\section{似然分数的近似：从理想到可行}

直接计算 $\nabla_{x_t} \log p(y | x_t)$ 是不可行的，因为 $p(y | x_t)$ 涉及从含噪中间数据 $x_t$ 到观测 $y$ 的复杂映射。我们需要一系列近似来使计算变得可行。

\subsection{关键困难：$x_t$ vs. $x_0$}

我们可以直接写出 $p(y | x_0)$ 的表达式——当 $\mathcal{A}$ 是线性函数时，$y | x_0 \sim \mathcal{N}(\mathcal{A} x_0, \sigma^2 I)$，这是一个简单的高斯分布。

但我们需要的是 $p(y | x_t)$，而不是 $p(y | x_0)$。问题在于：\textbf{如何建立 $x_t$ 和 $x_0$ 之间的联系？}

\begin{warningbox}{不能直接用 $p(y|x_0)$ 替代 $p(y|x_t)$}
最粗暴的近似是假设 $p(y|x_t) \approx p(y|x_0)$，即把中间噪声数据当作干净数据来计算。这是极差的近似——如果 $p(x_0|x_t)$ 真的是 delta 分布（即从 $x_t$ 可以确定性地恢复 $x_0$），那我们根本不需要迭代去噪过程了。
\end{warningbox}

\subsection{Tweedie 公式：连接 $x_t$ 与 $x_0$}

\textbf{Tweedie 公式}提供了从 $x_t$ 估计 $x_0$ 的均值的方法：

$$
\hat{x}_0 = \mathbb{E}[x_0 | x_t] = \frac{1}{\sqrt{\bar{\alpha}_t}} \left( x_t + (1 - \bar{\alpha}_t) \nabla_{x_t} \log p(x_t) \right)
$$

其中 $\nabla_{x_t} \log p(x_t)$ 正是预训练模型预测的分数（或等价地，通过预测噪声 $\epsilon_\theta(x_t, t)$ 得到）。

\begin{importantbox}{Tweedie 公式的核心作用}
Tweedie 公式告诉我们：给定 $x_t$，我们可以计算 $x_0$ 的\textbf{后验均值} $\hat{x}_0$（也称为去噪估计）。虽然这只是均值而非完整分布，但它为近似似然分数提供了关键的桥梁。
\end{importantbox}

\subsection{高斯近似：$p(x_0 | x_t)$ 的简化}

为了使计算可行，我们做出一个关键（但粗糙的）近似假设：

$$
p(x_0 | x_t) \approx \mathcal{N}(\hat{x}_0, r_t^2 I)
$$

即将 $x_0$ 在给定 $x_t$ 下的后验分布近似为一个以 Tweedie 估计 $\hat{x}_0$ 为均值、以某个标量方差 $r_t^2$ 为方差的高斯分布。

\begin{warningbox}{高斯近似的局限性}
这个近似是相当粗糙的。如果 $p(x_0|x_t)$ 真的是高斯分布，我们就不需要迭代去噪了——直接从 $x_t$ 采样 $x_0$ 即可。真实的 $p(x_0|x_t)$ 是一个远比高斯复杂的多模态分布。但在没有额外训练的情况下，这是我们能做到的最好近似之一。
\end{warningbox}

\subsection{本章小结}

通过 Tweedie 公式和高斯近似，我们建立了从 $x_t$ 到 $x_0$ 的桥梁，将不可直接计算的 $p(y|x_t)$ 转化为可以解析计算的高斯分布。接下来我们将基于此推导出具体的引导算法。


